\begin{textsample}{NEG Dim 2 – vem\_ed\_22 – Score 3.00 – vem\_ed\_22/t116.txt}  \label{ex:f2_neg_009}
PCI recebeu prêmio da Câmara de Comércio dos BRICS O Projeto Colaborativo Internacional (PCI / Cesu) sobre exportação de flores da África do Sul para o Mercosul, considerando os Objetivos de Desenvolvimento Sustentável (ODS) propostos na Agenda 2030 da Organização das Nações Unidas (ONU), recebeu prêmio da Câmara de Comércio e Indústria (CCI) dos BRICS - aliança composta por Brasil, Rússia, Índia, China e África do Sul, que recebe outras seis nações convidadas em 2024.
Realizado em 2022, o PCI envolveu professores e estudantes de Comércio Exterior das Fatecs Itapetininga, Indaiatuba e São Caetano do Sul e de Floricultura da Durban University of Technology (África do Sul).
João Guilherme Diniz, concluinte de Comércio Exterior na Fatec Itapetininga, ficou sabendo do prêmio e inscreveu o projeto.
A CCI reconheceu a contribuição para as parcerias estratégicas entre países dos BRICS e concedeu o prêmio de “Excelência Acadêmica”.
O prêmio “BRICS CCI Annual Recognition Awards” (BARA) destaca projetos que contribuam para o desenvolvimento sustentável dos países membros.
A professora Paula Granato esteve na cerimônia de premiação no hotel Le Méridien, em Nova Deli, em 19 de janeiro de 2024.
Aqui, o relato de Paula para o Instagram do Centro Paula Souza: https://shorturl.at/KLX37.
E a notícia publicada no site do Centro Paula Souza: https://www.cps.sp.gov.br/fatec-itapetininga-e-premiada-em-evento-na-india-no-ultimo-dia-19/

% matched lemmas: 
\end{textsample}
